\chapter{Introdução}
\label{1-introducao}

O presente trabalho está focado na Engenharia de Software, mais especificamente na área de segurança de software. Propõe-se o uso de Sistemas Multiagentes para viabilizar uma arquitetura P2P (\textit{Peer-to-Peer}) de agentes especialistas, semicentralizada via um Agente Auditor, para lidar com um portfólio de automações no contexto de práticas DevSecOps (Development, Security and Operations).

Este capítulo está organizado em seis seções: a Contextualização (seção~\ref{1-contextualizacao}), que situa o problema e o recorte de escopo desse problema a ser tratado no trabalho; a Questão de Pesquisa (seção~\ref{1-questao}), que orienta a investigação; a Justificativa (seção~\ref{1-justificativa}), que apresenta a relevância da proposta; os Objetivos (seção~\ref{1-objetivos}), divididos em objetivo geral e objetivos específicos; a Metodologia (seção~\ref{1-metodologia}), que apresenta uma visão geral da condução da pesquisa; e a Organização da Monografia (seção~\ref{1-organizacao}), que descreve o conteúdo dos demais capítulos.

\section{Contextualização}
\label{1-contextualizacao}

A segurança de software é primordial para as organizações que dependem de sistemas de informação e automações em suas atividades. Vulnerabilidades no código ou em seus componentes podem comprometer dados e serviços, o que torna necessária a identificação de falhas ao longo do ciclo de vida das soluções. Para integrar essas verificações ao desenvolvimento e à entrega de software, a abordagem DevSecOps articula desenvolvimento, segurança e operações \cite{owaspDevSecOps}. As esteiras de integração e entrega contínuas, conhecidas pela sigla CI/CD, organizam etapas como construção, teste, empacotamento e implantação, nas quais podem ser incorporados controles sobre o código e suas dependências \cite{chandramouli2024}.

Essa integração ainda exige esforço das pessoas. Segundo \citeonline{rajapakse2022}, a segurança é tradicionalmente tratada como requisito não funcional, verificada sobre uma versão já pronta, e atividades como a revisão de código demandam esforço manual substancial. As ferramentas de análise produzem evidências, mas a interpretação dessas evidências, continua a cargo dos responsáveis humanos, que decidem sobre os riscos encontrados.

Em organizações com um portfólio de projetos e automações, esse acompanhamento envolve analisar evidências provenientes de diferentes repositórios e considerar as particularidades de cada solução. Os relatórios técnicos precisam ser relacionados à finalidade da aplicação, às dependências utilizadas e às políticas pertinentes, para que os achados possam orientar uma recomendação, aqui denominada parecer. Quando esse acompanhamento se apoia apenas na análise humana dos resultados de ferramentas tradicionais de DevSecOps, o esforço analítico tende a ser significativo, conforme discutem \citeonline{zuev2025} ao tratarem do aumento da carga cognitiva associado à grande quantidade de informações providas via diversas ferramentas, bem como às restrições de tempo para avaliar riscos de segurança. O problema de interesse deste trabalho consiste em articular esses resultados, reunindo achados de diferentes dimensões de análise e interpretando-os de maneira que o responsável humano possa examinar a fundamentação do parecer produzido.

A diversidade das formas de criação das automações também compõe esse contexto de análise. Plataformas \textit{low-code/no-code} (LCNC) permitem construir aplicações e automações com pouco ou nenhum código, e favorecem o \textit{citizen development}, isto é, a participação de usuários de negócio na criação de soluções. A adoção dessas práticas envolve benefícios e desafios organizacionais e técnicos, incluindo questões de segurança e governança \cite{ajimati2025}. Quando as soluções incorporam inteligência artificial fora das estruturas de governança aprovadas pela organização, o fenômeno denominado \textit{Shadow AI} envolve questões de privacidade dos dados e de responsabilidade sobre o uso da tecnologia \cite{silic2025}. Essa variedade de autores, práticas e procedimentos de revisão pode ampliar a heterogeneidade do portfólio e tornar mais relevante o acompanhamento contínuo dos repositórios, com apoio de práticas de DevSecOps e de auditorias técnicas.

Diante do esforço que esse acompanhamento exige das pessoas, alternativas de soluções são desejadas. Nesse contexto, os Sistemas Multiagentes (SMA) oferecem uma possibilidade de organizar o apoio à análise por meio de agentes com responsabilidades especializadas \cite{wooldridge2009}. A comunicação estruturada permite solicitar informações e compartilhar evidências \cite{lazarin2025}, enquanto a definição explícita de papéis e procedimentos decisórios fornece uma base para organizar a colaboração humano-agente \cite{ait2025}.

O paradigma de Sistemas Multiagentes é útil para lidar com situações que incorrem em descentralização, partindo para uma solução distribuída \cite{wang2026}. Nesse sentido, uma solução orientada a Sistema Multiagentes demanda uma arquitetura compatível com essa distribuição, sendo comumente utilizada a Arquitetura \textit{Peer-to-Peer} (P2P) com algum mecanismo de controle semi-centralizado de preferência \cite{moro2003,flores2005}. No contexto em estudo neste trabalho, por meio da comunicação direta P2P, os agentes poderão solicitar esclarecimentos e examinar informações de outras dimensões do problema, antes de encaminhar suas conclusões ao Agente Auditor, sendo esse responsável pelo controle semi-centralizado. A ideia é que o Agente Auditor consolide, por regras explícitas, um parecer que relacione a recomendação centrada em segurança às evidências consideradas pelos demais agentes, identificando situações que demandem encaminhamento ao responsável humano. 

Dada a abrangência de cenários que demandam soluções de segurança como foco, o recorte desse domínio, que limita o escopo de atuação deste trabalho, insere-se em prover uma solução piloto, concebida como prova de conceito (PoC), de uma arquitetura multiagentes P2P semi-centralizada, visando apoiar a auditoria de repositórios do portfólio de automações de uma empresa de porte grande, com 120 funcionários contratados, e um portfólio de aproximadamente 22 projetos, atuando na região do DF.

Ao longo da realização desse projeto piloto, será conduzido um acompanhamento de viés observacional, documentado com registro dos achados, das interações entre agentes especialistas, das justificativas e dos pareceres. Os achados documentados serão confrontados com fontes adicionais (ex. literatura e boas práticas adotadas na comunidade de segurança de software), permitindo perceber sobre a pertinência da solução ao contexto apresentado. Essas fontes constam mais bem acordadas no ~\ref{cap-referencial-teorico}. A análise será complementada por meio de uma apreciação humana quanto à clareza, à coerência e à fundamentação dos pareceres. Pretende-se que essa parte descritiva e documental do projeto seja centrada nas contribuições e nos limites observados, sem estimar efeitos causais.

\section{Questão de Pesquisa}
\label{1-questao}

Para delimitar o escopo e guiar a investigação, formula-se a seguinte questão de pesquisa.

\begin{quote}
\textit{Como apoiar a auditoria de repositórios de um portfólio de automações no contexto de práticas DevSecOps usando um sistema de agentes de software especialistas, estruturado via uma arquitetura P2P semi-centralizada em um Agente Auditor?}
\end{quote}


\section{Justificativa}
\label{1-justificativa}

O trabalho justifica-se pela dificuldade de acompanhar portfólios de projetos e automações quando a análise dos resultados de ferramentas tradicionais recai exclusivamente sobre pessoas, o que torna necessário organizar as evidências técnicas em um processo de apoio à decisão, conforme sugerido em \citeonline{paulo2026}. A autora Sofia Marinho Paulo destaca que a "documentação de \textit{releases} de software em ambientes DevSecOps empresariais é um processo crítico que suporta decisões de aprovação, auditoria e rastreabilidade de alterações", e enfatiza que essa prática é realizada maioritariamente de forma manual, envolvendo a consulta de vários insumos gerados (ex. \textit{commits} e \textit{pull requests}), além da compilação da documentação. Isso torna o processo oneroso, inconsistente e de difícil auditoria. Afirma-se ainda, com base em uma revisão sistemática da literatura, que praticamente não existem suportes de apoio que considerem um fluxo integrado, envolvendo geração automática validada e rastreabilidade verificável, orientado a contextos DevSecOps empresariais. 

\citeonline{zuev2025} oferecem uma referência para transformar resultados de análise em informação para decisão. \citeonline{ait2025} e \citeonline{yan2025} contribuem para discutir políticas explícitas e raciocínio normativo. \citeonline{mosciatti2026} e \citeonline{sowan2026} oferecem elementos para delimitar permissões, registros e apreciação humana em seus respectivos domínios. A integração desses elementos na auditoria de repositórios constitui o recorte do projeto. Ressalta-se ainda que, mesmo ocorrendo menção à falta de soluções integradas nesse contexto \citeonline{paulo2026}, o trabalho proposto não pressupõe análises sobre a eficácia já demonstrada ou ausência de soluções equivalentes na literatura.

A arquitetura P2P semi-centralizada permitirá que um agente especialista solicite esclarecimentos sobre evidências de outras nuances da situação em análise, pois a interpretação de um achado pode depender dessas evidências. Um alerta de segurança, por exemplo, pode ser compreendido de modo diferente, quando um agente de qualidade informa que o mesmo módulo concentra alta complexidade, ou ainda, quando um agente requisitos indica a finalidade e a criticidade do módulo, respectivamente, como central e alta.

O projeto piloto abrange repositórios do portfólio de automações da Empresa A. Esse portfólio reúne automações de origens diversas, e a presença de automações desenvolvidas por usuários de negócio, de soluções construídas com plataformas \textit{low-code/no-code} e de usos de inteligência artificial fora das estruturas de governança aprovadas, isto é, de situações de \textit{Shadow AI}. Em pesquisa com 88 profissionais, \citeonline{canedo2026} registram maturidade desigual em DevSecOps e, nas respostas abertas, a preocupação com o uso de inteligência artificial fora do controle institucional. A caracterização do contexto institucional será apresentada com base nas informações disponíveis e autorizadas para a pesquisa, preservando o sigilo da organização. O uso dos repositórios depende de autorização institucional e de condições de divulgação, o que já está sendo negociado pelo autor dessa proposta junto à Empresa A. Até o momento, há indicação de abertura suficiente para a condução do trabalho.

Do ponto de vista acadêmico, pretende-se documentar o projeto como um todo, com ênfase na teoria, nas tecnologias, na metodologia, nos insumos desenvolvidos e nas análises. Esses registros poderão ser consultados por outros interessados além de subsidiarem o refinamento futuro do projeto piloto. Por si só, essa parte documental do projeto já representa uma contribuição da proposta, também sendo destacada como justificativa para a realização da mesma. A análise permanecerá delimitada aos repositórios autorizados pela Empresa A, bem como às condições efetivamente observáveis.

\section{Objetivos}
\label{1-objetivos}

Para orientar o desenvolvimento do projeto piloto e a análise de sua aplicação em uma empresa real, foram definidos o objetivo geral e alguns objetivos específicos, apresentados a seguir.

\subsection{Objetivo Geral}
\label{1-objetivos-gerais}

Apoio à auditoria de repositórios de um portfólio de automações em uma empresa real, usando um sistema multiagentes especializado em segurança de software (práticas DevSecOps) e estruturado como uma arquitetura P2P semi-centralizada.

\subsection{Objetivos Específicos}
\label{1-objetivos-especificos}

Os objetivos específicos, em conjunto, buscam o cumprimento do objetivo geral, abrangendo a fundamentação teórica, a modelagem, a implementação e a avaliação da proposta.

\begin{itemize}
    \item \textbf{Objetivo Específico 1:} Levantamento teórico e tecnológico sobre Sistemas Multiagentes, arquitetura P2P semi-centralizada e segurança de software, com foco em governança de software, auditoria e práticas DevSecOps;
    \item \textbf{Objetivo Específico 2:} Levantamento de trabalhos relacionados a sistemas multiagentes, comunicação entre agentes e apoio à decisão em auditoria, com ênfase em segurança de software(práticas DevSecOps);
    \item \textbf{Objetivo Específico 3:} Modelagem das responsabilidades dos agentes, do fluxo de análise e dos mecanismos de comunicação na arquitetura P2P semi-centralizada;
    \item \textbf{Objetivo Específico 4:} Implementação do projeto piloto, considerando os insumos acordados com os levantamentos anteriores, bem como a modelagem arquitetural dos agentes;
    \item \textbf{Objetivo Específico 5:} Rastreabilidade das interações entre agentes especialistas, incluindo o registro das evidências compartilhadas e das condições para se manter o encaminhamento sobre segurança; ou ajustá-lo em função das informações advindas das interações;
    \item \textbf{Objetivo Específico 6:} Incorporação do Agente Auditor à solução para consolidar os pareceres, fundamentando-se em insumos de Governança de Software em Segurança (ex. controles técnicos centrados nos objetivos de negócio da empresa e em conformidade às leis e normas regulatórias cabíveis), e identificando até mesmo a necessidade complementar de análise humana;
    \item \textbf{Objetivo Específico 7:} Análise dos resultados por meio de um estudo observacional e descritivo do projeto piloto, considerando a rastreabilidade das evidências e das interações entre agentes, além da apreciação por especialistas dos pareceres gerados, e
    \item \textbf{Objetivo Específico 8:} Documentação dos achados de pesquisa, de desenvolvimento e de análise, visando divulgar a solução, bem como aperfeiçoar essa solução em trabalhos futuros.
\end{itemize}

\section{Metodologia}
\label{1-metodologia}

De acordo com \cite{gil2002}, a pesquisa pode ser compreendida em termos de natureza, objetivos, abordagem e procedimentos. Em termos de natureza, trata-se de uma pesquisa aplicada, sendo a solução proposta para o portfólio de uma empresa real. Já quanto aos objetivos, a pesquisa é exploratória e descritiva, demandando lidar com um portfólio de projetos reais, em aspectos de segurança, e buscando uma solução orientada a um paradigma emergente. Tanto o portfólio, quanto os tópicos de segurança de software e sistemas multiagentes são insumos em estudo por parte do autor, o que torna a proposta ainda mais desafiadora e exploratória. Além disso, pretende-se documentar e explicar cada nuance de relevância do projeto piloto, o que sugere uma pesquisa baseada em rastros descritivos sobre cada ponto observado. Quanto à abordagem, a pesquisa é predominantemente qualitativa sendo analisada, principalmente, com base na percepção de especialistas em segurança, atuantes no mercado. Mas, não é descartado o uso de métricas de cunho quantitativo, quando ocorrer a necessidade de evidenciar aspectos de forma mais concreta e clara, por exemplo, sobre o desempenho dos agentes, o número de interações, ou mesmo sobre a taxa de acerto dos agentes, quando os pareceres gerados forem avaliados pelos especialistas humanos. Por fim, os principais procedimentos serão: (i) pesquisa bibliográfica, para embasamento dos tópicos em estudo na literatura; (ii) prova de conceito, evidenciando que o projeto piloto não pressupõe gerar uma solução plena, implantada na empresa, mas sim uma prova de conceito que sugere um caminho de solução para apoio à auditoria de repositórios de um portfólio de automações no contexto de práticas DevSecOps, e (iii) estudo observacional e descritivo, buscando manter o projeto piloto documentado para consultas e trabalhos futuros.

Dianto do exposto, a solução será examinada por meio de um estudo observacional e descritivo, em versões identificadas de repositórios autorizados pela Empresa A. Pretende-se documentar os achados provenientes do uso de ferramentas e práticas DevSecOps; de interações entre agentes e pareceres produzidos. A análise focará na coerência desses achados com suas fontes, leis e normas. O planejamento da coleta considerará, por adaptação ao domínio de software, as contribuições metodológicas de \citeonline{zangirolami2018} e de \citeonline{morshed2009}, além das diretrizes de planejamento da coleta de dados e de relato de estudos empíricos em engenharia de software apresentadas por \citeonline{runeson2009}.

Como avaliação complementar, será realizada uma apreciação humana, captando a percepção de especialistas em segurança sobre a clareza, a coerência e a fundamentação dos pareceres gerados pelo sistema de multiagentes. A classificação detalhada da pesquisa; a pesquisa bibliográfica, os critérios de seleção dos repositórios e dos avaliadores; os métodos que orientam o desenvolvimento (i.r. uma combinação de Scrum Solo(\citeonline{pagotto2016}) e Kanban (\citeonline{boeg2010})); os instrumentos de coleto e os procedimentos de análise serão apresentados no Capítulo~\ref{cap-metodologia}.

\section{Organização da Monografia}
\label{1-organizacao}

A monografia está organizada em seis capítulos. Após esta introdução, os demais capítulos apresentam a fundamentação, as tecnologias, os procedimentos metodológicos, a proposta e as considerações finais sobre a primeira etapa do trabalho, conforme descrito a seguir:

\begin{itemize}
    \item \textbf{Capítulo~\ref{cap-referencial-teorico} -- Referencial Teórico:} apresenta os conceitos e os trabalhos relacionados que fundamentam a pesquisa, abrangendo segurança de software(mais especificamente, auditoria de software e práticas DevSecOps) e apoio à decisão, sistemas multiagentes, interpretação documental apoiada por modelos de linguagem, comunicação entre agentes especialistas, políticas de governança, explicabilidade, rastreabilidade e revisão humana;

    \item \textbf{Capítulo~\ref{cap-suporte-tecnologico} -- Suporte Tecnológico:} descreve as tecnologias para o projeto piloto, a infraestrutura de comunicação entre agentes, os recursos de interpretação documental e as ferramentas de análise de qualidade, segurança e dependências. Também apresenta os recursos de apoio geral, apoio à modelagem e ao registro das evidências;

    \item \textbf{Capítulo~\ref{cap-metodologia} -- Metodologia:} detalha a classificação da pesquisa, o levantamento bibliográfico, o método de desenvolvimento da solução piloto e o planejamento da avaliação observacional e descritiva. Explicita os procedimentos de seleção dos repositórios, coleta de evidências, apreciação humana e análise dos registros, além dos processos e dos cronogramas de atividades para ambas as etapas do trabalho;

    \item \textbf{Capítulo~\ref{cap-proposta} -- Proposta:} apresenta a Empresa A, o perfil do portfólio de repositórios da mesma, o escopo e os requisitos da solução, a arquitetura multi-agente e as responsabilidades dos participantes. Detalha a comunicação P2P entre agentes especialistas, a consolidação de pareceres pelo Auditor e os mecanismos previstos para conservar evidências, justificativas e situações que demandem revisão humana, e

    \item \textbf{Capítulo~6 -- Considerações Finais:} retoma os objetivos e sintetiza o estágio de desenvolvimento alcançado na primeira etapa do trabalho, apresentando as limitações identificadas e os próximos passos para a continuidade da implementação e da avaliação da solução, além de esclarecimentos sobre o uso de IA Generativa no trabalho.
\end{itemize}

Ao final, são apresentadas as referências bibliográficas e demais insumos gerados ao longo do trabalho em apêndices (ex. o instrumento proposta para o registro das auditorias)